\begin{resumen}
La proliferación de las redes sociales basadas en vídeos cortos, como \textit{TikTok}, ha traído consigo un aumento exponencial del discurso de odio y la misoginia en entornos digitales. La detección automática de este contenido es un desafío crítico, ya que el sexismo rara vez se manifiesta de forma explícita, ocultándose a menudo tras el sarcasmo, los estereotipos o el contexto visual de la escena. Este Trabajo de Fin de Máster aborda el problema proponiendo un sistema de clasificación multimodal avanzado, diseñado y evaluado bajo el marco del reto oficial \textit{EXIST 2025}.

A nivel metodológico, la investigación amplía el enfoque clásico unimodal basado únicamente en el texto. Se extraen y procesan las tres dimensiones fundamentales del contenido multimedia: el discurso verbal (transcripción), la disposición espacial (imágenes estáticas) y el movimiento a lo largo del tiempo (secuencia de vídeo). Para ello, se entrenan arquitecturas de estado del arte de forma independiente, aplicando técnicas de \textit{Fine-Tuning} y cuantización (\textit{QLoRA}) sobre modelos discriminativos (\textit{RoBERTa}), modelos de lenguaje generativos masivos (\textit{Mistral-7B}), y arquitecturas de visión (\textit{ConvNeXt} y \textit{TimeSFormer}). Las predicciones de los mejores modelos expertos se unifican posteriormente mediante una estrategia de fusión tardía (\textit{Late Fusion}). Además, la investigación explora el paradigma del Perspectivismo (\textit{Learning with Disagreement}), analizando cómo la identidad de género del anotador podría influir en el entrenamiento de los modelos. 

Los resultados empíricos preliminares apuntan a una mejora de la arquitectura híbrida frente a los enfoques aislados, logrando un \textit{F1-Score} de 0.7261 sobre un subconjunto de prueba local utilizado para la calibración. Al ser evaluado de forma independiente frente al conjunto ciego del reto, el sistema alcanzó un \textit{F1-Score} final de 0.6311, respaldando la viabilidad de la propuesta multimodal. El análisis multietiqueta sugiere que el texto destaca como una modalidad efectiva para detectar sesgos ideológicos, mientras que la información visual y temporal parece aportar un contexto útil para fenómenos complejos como la cosificación. No obstante, se identifican limitaciones relevantes ligadas al fuerte desbalance de clases, el sobreajuste local y el ruido derivado de la extracción automática de texto, lo que subraya el carácter acotado de estos resultados. Finalmente, los experimentos de perspectivismo muestran variaciones en las métricas de recuperación (\textit{Recall}) según el género del anotador, sugiriendo que la demografía influye en la evaluación, si bien estos hallazgos deben interpretarse con cautela debido a los sesgos y dimensiones de las distribuciones de los datos.

\noindent\textbf{Palabras clave}: Detección de Sexismo, Aprendizaje Profundo, Multimodalidad, Análisis del Discurso de Odio, Perspectivismo.
\end{resumen}

\clearpage

\begin{abstract}
The proliferation of short-video social networks, such as \textit{TikTok}, has led to an exponential increase in hate speech and misogyny in digital environments. The automatic detection of this content is a critical challenge, as sexism rarely manifests explicitly, often hiding behind sarcasm, stereotypes, or the visual context of the scene. This Master's Thesis addresses the problem by proposing an advanced multimodal classification system, designed and evaluated under the framework of the official \textit{EXIST 2025} challenge. 

Methodologically, the research expands upon the classic text-only unimodal approach. The three fundamental dimensions of multimedia content are extracted and processed: verbal discourse (transcription), spatial arrangement (static images), and movement over time (video sequences). To this end, state-of-the-art architectures are trained independently, applying \textit{Fine-Tuning} and quantization (\textit{QLoRA}) techniques on discriminative models (\textit{RoBERTa}), massive generative language models (\textit{Mistral-7B}), and vision architectures (\textit{ConvNeXt} and \textit{TimeSFormer}). The predictions of the best expert models are subsequently unified using a \textit{Late Fusion} strategy. Furthermore, the research explores the Perspectivism paradigm (\textit{Learning with Disagreement}), analyzing how the annotator's gender identity might influence model training. 

Preliminary empirical results point to an improvement of the hybrid architecture over isolated approaches, achieving an \textit{F1-Score} of 0.7261 on a local test subset used for calibration. When evaluated independently against the challenge's blind test set, the system achieved a final \textit{F1-Score} of 0.6311, supporting the viability of the multimodal proposal. Multi-label analysis suggests that text stands out as an effective modality for detecting ideological biases, while visual and temporal information seems to provide useful context for complex phenomena such as objectification. However, relevant limitations are identified regarding severe class imbalance, local overfitting, and noise from automatic text extraction, highlighting the bounded nature of these results. Finally, perspectivism experiments show variations in retrieval metrics (\textit{Recall}) according to the annotator's gender, suggesting that demographics influence evaluation, although these findings must be interpreted with caution due to biases and data distribution sizes.

\bigskip
\noindent\textbf{Keywords}: Sexism Detection, Deep Learning, Multimodality, Hate Speech Analysis, Perspectivism. 
\end{abstract}
